\clearpage
\begin{landscape}
\begingroup
\fontsize{9.5}{12}\selectfont
\fancyhead[L]{\small NACHLASS. STÖRUNGEN DER PALLAS}
\renewcommand{\headrulewidth}{0pt}
\makeatletter\tagsleft@true\let\veqno\@@leqno\makeatother
\setlength{\lineskip}{4pt}\setlength{\lineskiplimit}{4pt}
\newcommand{\VthreeTerminalPair}{\mathord{\text{\makebox[.14em][l]{\raisebox{.75ex}{\rule{.055em}{.16ex}}\kern-.055em\raisebox{-.18ex}{,}}}}}
\newcommand{\VthreeRaisedPoint}{\mathord{\text{\raisebox{.5ex}{.}}}}
\newcommand{\VthreeLowDash}{\mathord{\text{\makebox[0pt][l]{\kern-.06em\raisebox{-.32ex}{\rule{.15em}{.045ex}}}}}}
\newcommand{\VthreeBrokenMark}{\mathord{\text{\makebox[.1em][l]{\raisebox{.8ex}{\rule{.055em}{.22ex}}\kern-.055em\raisebox{-.1ex}{\rule{.08em}{.5ex}}}}}}
\newcommand{\VthreeLowMark}{\mathord{\text{\makebox[0pt][l]{\kern-.1em\raisebox{-.65ex}{,}}}}}
\newcommand{\VthreeVerticalMark}{\mathord{\text{\makebox[.1em][l]{\raisebox{-.1ex}{\rule{.08em}{1.1ex}}}}}}
\SourceLeaf{502}
\pdfbookmark[2]{Allgemeine Störungen, Erste Rechnung, [3]}{pallas-V-fragment3}
\begin{center}[3.]\end{center}
\par\noindent\hspace*{\parindent}[Die im vorstehenden Artikel abgeleiteten Formeln sollen nun für die bei der Berechnung der Pallasstörungen angewandten Fälle specialisirt werden, nemlich für \(\mu=12,24,48\), wobei fast überflüssig ist zu bemerken, dass man alle vorkommenden Winkel, welche grösser als \(90^0\) sind, durch solche aus dem ersten Quadranten sofort ersetzt.
\par\noindent\hspace*{\parindent}Es soll zunächst die Specialisirung für \(\mu=48\) ausgeführt werden. Hier wird das grösste vorkommende Vielfache von \(a\) gleich \(11a=82^0 30'\) sein.
\par\noindent\hspace*{\parindent}Man erhält mit Hülfe der Formeln 12 des vorigen Artikels:
\[
\tag*{1)}
\begin{aligned}
24\alpha_4&=\tfrac12 M_0+M_1\cos4a+M_2\cos8a\\
24\beta_4&=N_1\sin4a+N_2\sin8a+\tfrac12 N_3\\[4pt]
24\alpha_{12}&=\tfrac12 M_0-M_2\\
24\beta_{12}&=N_1-\tfrac12 N_3\\[4pt]
24\alpha_{20}&=\tfrac12 M_0-M_1\cos4a+M_2\cos8a\\
24\beta_{20}&=N_1\sin4a-N_2\sin8a+\tfrac12 N_3.
\end{aligned}
\]
\par\noindent\hspace*{\parindent}Ferner aus den Formeln 14:
\[
\tag*{2)}
\begin{aligned}
24\alpha_8&=\tfrac12(P_0-P_3)+(P_1-P_2)\cos8a\\
24\beta_8&=(Q_1+Q_2)\sin8a\VthreeTerminalPair\\[4pt]
24\alpha_{16}&=\tfrac12(P_0+P_3)-(P_1+P_2)\cos8a\\
24\beta_{16}&=(Q_1-Q_2)\sin8a,
\end{aligned}
\]
denen man die beiden:
\[
\begin{aligned}
48\alpha_0&=\tfrac12(P_0+P_3)+(P_1+P_2)\\
48\alpha_{24}&=\tfrac12(P_0-P_3)-(P_1-P_2)]
\end{aligned}
\]
beifügen kann.
\par\noindent\hspace*{\parindent}Ferner ergibt sich aus 8:
\[
\tag*{3)}
\begin{aligned}
24\alpha_2&=\tfrac12 H_0+H_1\cos2a+H_2\cos4a+H_3\cos6a+H_4\cos8a+H_5\cos10a\\
24\beta_2&=\mathrm{J}_1\sin2a+\mathrm{J}_2\sin4a+\mathrm{J}_3\sin6a+\mathrm{J}_4\sin8a+\mathrm{J}_5\sin10a+\tfrac12\mathrm{J}_6\\[4pt]
24\alpha_6&=\tfrac12 H_0-H_4+(H_1-H_3-H_5)\cos6a\\
24\beta_6&=\mathrm{J}_2-\tfrac12\mathrm{J}_6+(\mathrm{J}_1+\mathrm{J}_3-\mathrm{J}_5)\sin6a\\[4pt]
24\alpha_{10}&=\tfrac12 H_0+H_1\cos10a-H_2\cos4a-H_3\cos6a+H_4\cos8a+H_5\cos2a\\
24\beta_{10}&=\mathrm{J}_1\sin10a+\mathrm{J}_2\sin4a-\mathrm{J}_3\sin6a-\mathrm{J}_4\sin8a+\mathrm{J}_5\sin2a+\tfrac12\mathrm{J}_6\\[4pt]
24\alpha_{14}&=\tfrac12 H_0-H_1\cos10a-H_2\cos4a+H_3\cos6a+H_4\cos8a-H_5\cos2a\\
24\beta_{14}&=\mathrm{J}_1\sin10a-\mathrm{J}_2\sin4a-\mathrm{J}_3\sin6a+\mathrm{J}_4\sin8a+\mathrm{J}_5\sin2a-\tfrac12\mathrm{J}_6\\[4pt]
24\alpha_{18}&=\tfrac12 H_0-H_4-(H_1-H_3-H_5)\cos6a\\
24\beta_{18}&=-\mathrm{J}_2+\tfrac12\mathrm{J}_6+(\mathrm{J}_1+\mathrm{J}_3-\mathrm{J}_5)\sin6a\\[4pt]
24\alpha_{22}&=\tfrac12 H_0-H_1\cos2a+H_2\cos4a-H_3\cos6a+H_4\cos8a-H_5\cos10a\\
24\beta_{22}&=\mathrm{J}_1\sin2a-\mathrm{J}_2\sin4a+\mathrm{J}_3\sin6a-\mathrm{J}_4\sin8a+\mathrm{J}_5\sin10a-\tfrac12\mathrm{J}_6.
\end{aligned}
\]
\clearpage\SourceLeaf{503}
\par\noindent\hspace*{\parindent}Endlich aus 6:
\[
\tag*{4)}
\begin{aligned}
12(\alpha_1+\alpha_{23})&=\tfrac12 D_0+D_2\cos2a+D_4\cos4a+D_6\cos6a+D_8\cos8a+D_{10}\cos10a\\
12(\beta_1-\beta_{23})&=E_2\sin2a+E_4\sin4a+E_6\sin6a+E_8\sin8a+E_{10}\sin10a+\tfrac12 E_{12}\\[4pt]
12(\alpha_3+\alpha_{21})&=\tfrac12 D_0-D_8+(D_2-D_6-D_{10})\cos6a\\
12(\beta_3-\beta_{21})&=E_4-\tfrac12 E_{12}+(E_2+E_6-E_{10})\sin6a\\[4pt]
12(\alpha_5+\alpha_{19})&=\tfrac12 D_0+D_2\cos10a-D_4\cos4a-D_6\cos6a+D_8\cos8a+D_{10}\cos2a\\
12(\beta_5-\beta_{19})&=E_2\sin10a+E_4\sin4a-E_6\sin6a-E_8\sin8a+E_{10}\sin2a+\tfrac12 E_{12}\\[4pt]
12(\alpha_7+\alpha_{17})&=\tfrac12 D_0-D_2\cos10a-D_4\cos4a+D_6\cos6a+D_8\cos8a-D_{10}\cos2a\\
12(\beta_7-\beta_{17})&=E_2\sin10a-E_4\sin4a-E_6\sin6a+E_8\sin8a+E_{10}\sin2a-\tfrac12 E_{12}\\[4pt]
12(\alpha_9+\alpha_{15})&=\tfrac12 D_0-D_8-(D_2-D_6-D_{10})\cos6a\VthreeTerminalPair\\
12(\beta_9-\beta_{15})&=-E_4+\tfrac12 E_{12}+(E_2+E_6-E_{10})\sin6a\\[4pt]
12(\alpha_{11}+\alpha_{13})&=\tfrac12 D_0-D_2\cos2a+D_4\cos4a-D_6\cos6a+D_8\cos8a-D_{10}\cos10a\\
12(\beta_{11}-\beta_{13})&=E_2\sin2a-E_4\sin4a+E_6\sin6a-E_8\sin8a+E_{10}\sin10a-\tfrac12 E_{12}
\end{aligned}
\]
und die übrigen, indem man theilweise nochmals ihre Summe resp. Differenz bildet:
\[
\tag*{5)}
\begin{aligned}
12(\alpha_1+\alpha_7-\alpha_{17}-\alpha_{23})&=\VthreeLowDash(D_1+D_7)(\cos a+\cos7a)+(D_5-D_{11})(\cos5a-\cos11a)\\
12(\beta_1-\beta_7-\beta_{17}+\beta_{23})&=(E_1-E_7)(\sin a-\sin7a)+(E_5+E_{11})(\sin5a+\sin11a)\\[4pt]
12(\alpha_1-\alpha_7+\alpha_{17}-\alpha_{23})&=(D_1-D_7)(\cos a-\cos7a)+(D_5+D_{11})(\cos5a+\cos11a)+2D_3\cos3a+2D_9\cos9a\\
12(\beta_1+\beta_7+\beta_{17}+\beta_{23})&=(E_1+E_7)(\sin a+\sin7a)+(E_5-E_{11})(\sin5a-\sin11a)+2E_3\sin3a+2E_9\sin9a\\[4pt]
12(\alpha_3-\alpha_{21})&=(D_1-D_7-D_9)\cos3a+(D_3-D_5-D_{11})\cos9a\\
\VthreeBrokenMark12(\beta_3+\beta_{21})\VthreeLowMark&=(E_1+E_7-E_9)\sin3a+(E_3+E_5-E_{11})\sin9a\\[4pt]
12(\alpha_5+\alpha_{11}-\alpha_{13}-\alpha_{19})&=(D_1-D_7)(\cos5a+\cos11a)-(D_5+D_{11})(\cos a-\cos7a)-2D_3\cos9a+2D_9\cos3a\\
12(\beta_5-\beta_{11}-\beta_{13}+\beta_{19})&=(E_1+E_7)(\sin5a-\sin11a)-(E_5-E_{11})(\sin a+\sin7a)+2E_3\sin9a-2E_9\sin3a\\[4pt]
12(\alpha_5-\alpha_{11}+\alpha_{13}-\alpha_{19})&=(D_1+D_7)(\cos5a-\cos11a)-(D_5-D_{11})(\cos a+\cos7a)\\
12(\beta_5+\beta_{11}+\beta_{13}+\beta_{19})&=(E_1-E_7)(\sin5a+\sin11a)-(E_5+E_{11})(\sin a-\sin7a)\\[4pt]
12(\alpha_9-\alpha_{15})&=-(D_3-D_5-D_{11})\cos3a+(D_1-D_7-D_9)\cos9a\\
\VthreeVerticalMark12(\beta_9+\beta_{15})&=-(E_3+E_5-E_{11})\sin3a+(E_1+E_7-E_9)\sin9a.
\end{aligned}
\]
\par\noindent\hspace*{\parindent}In allen diesen Gleichungen ist \(a=7^0 30'\), welcher Werth aber der grössern Übersichtlichkeit wegen nicht eingesetzt worden ist. Es lassen sich hienach sämmtliche Coefficienten berechnen; die Gleichungen können aber zur numerischen Rechnung noch, wie folgt, vereinfacht werden; es ist nemlich \(\cos2a=\tfrac14(\BareRoot{6}+\BareRoot{2})\), \(\cos4a=\tfrac12\BareRoot{3}\), \(\cos6a=\tfrac12\BareRoot{2}\), \(\cos8a=\tfrac12\), \(\cos10a=\tfrac14(\BareRoot{6}-\BareRoot{2})\); setzt man diese Werthe und die entsprechenden für die Sinus ein, so wird aus 1:
\[
\tag*{I.}
\begin{aligned}
\alpha_{12}&=\frac{M_0-2M_2}{48}&\qquad\beta_{12}&=\frac{2N_1-N_3}{48}\\[6pt]
\alpha_4&=\frac{M_0+M_2}{48}+\frac{M_1}{\BareRoot{768}}&\beta_4&=\frac{N_1+N_3}{48}+\frac{N_2}{\BareRoot{768}}\\[6pt]
\alpha_{20}&=\frac{M_0+M_2}{48}-\frac{M_1}{\BareRoot{768}}\VthreeRaisedPoint&\beta_{20}&=\frac{N_1+N_3}{48}-\frac{N_2}{\BareRoot{768}}\VthreeRaisedPoint
\end{aligned}
\]
\clearpage\SourceLeaf{504}
\par\noindent\hspace*{\parindent}Ferner aus 2:
\[
\tag*{II.}
\begin{aligned}
\alpha^0&=\frac{P_0+P_3}{96}+\frac{P_1+P_2}{48}&&\\[6pt]
\alpha_{16}&=\frac{P_0+P_3}{48}-\frac{P_1+P_2}{48}&\qquad\beta_{16}&=\frac{Q_1-Q_2}{\BareRoot{768}}\\[6pt]
\alpha_8&=\frac{P_0-P_3}{48}+\frac{P_1-P_2}{48}&\beta_8&=\frac{Q_1+Q_2}{\BareRoot{768}}\\[6pt]
\alpha_{24}&=\frac{P_0-P_3}{96}-\frac{P_1-P_2}{48}\VthreeRaisedPoint&&
\end{aligned}
\]
\par\noindent\hspace*{\parindent}Ferner aus 3:
\[
\tag*{III.}
\begin{aligned}
\alpha_6&=\frac{H_0-2H_4}{48}+\frac{H_1-H_3-H_5}{\BareRoot{1152}}&\qquad\beta_6&=\frac{2\mathrm{J}_2-\mathrm{J}_6}{48}+\frac{\mathrm{J}_1+\mathrm{J}_3-\mathrm{J}_5}{\BareRoot{1152}}\\[6pt]
\alpha_{18}&=\frac{H_0-2H_4}{48}-\frac{H_1-H_3-H_5}{\BareRoot{1152}}&\beta_{18}&=-\frac{2\mathrm{J}_2-\mathrm{J}_6}{48}+\frac{\mathrm{J}_1+\mathrm{J}_3-\mathrm{J}_5}{\BareRoot{1152}}\\[6pt]
\alpha_2&=\frac{H_0+H_4}{48}+\frac{H_2}{\BareRoot{768}}+\frac{H_1+H_5}{\BareRoot{1536}}+\frac{H_1-H_5+2H_3}{\BareRoot{4608}}&\beta_2&=\frac{\mathrm{J}_2+\mathrm{J}_6}{48}+\frac{\mathrm{J}_4}{\BareRoot{768}}+\frac{\mathrm{J}_1+\mathrm{J}_5}{\BareRoot{1536}}-\frac{\mathrm{J}_1-\mathrm{J}_5-2\mathrm{J}_3}{\BareRoot{4608}}\\[6pt]
\alpha_{10}&=\frac{H_0+H_4}{48}-\frac{H_2}{\BareRoot{768}}+\frac{H_1+H_5}{\BareRoot{1536}}-\frac{H_1-H_5+2H_3}{\BareRoot{4608}}&\beta_{10}&=\frac{\mathrm{J}_2+\mathrm{J}_6}{48}-\frac{\mathrm{J}_4}{\BareRoot{768}}+\frac{\mathrm{J}_1+\mathrm{J}_5}{\BareRoot{1536}}+\frac{\mathrm{J}_1-\mathrm{J}_5-2\mathrm{J}_3}{\BareRoot{4608}}\\[6pt]
\alpha_{14}&=\frac{H_0+H_4}{48}-\frac{H_2}{\BareRoot{768}}-\frac{H_1+H_5}{\BareRoot{1536}}+\frac{H_1-H_5+2H_3}{\BareRoot{4608}}&\beta_{14}&=-\frac{\mathrm{J}_2+\mathrm{J}_6}{48}+\frac{\mathrm{J}_4}{\BareRoot{768}}+\frac{\mathrm{J}_1+\mathrm{J}_5}{\BareRoot{1536}}+\frac{\mathrm{J}_1-\mathrm{J}_5-2\mathrm{J}_3}{\BareRoot{4608}}\\[6pt]
\alpha_{22}&=\frac{H_0+H_4}{48}+\frac{H_2}{\BareRoot{768}}-\frac{H_1+H_5}{\BareRoot{1536}}-\frac{H_1-H_5+2H_3}{\BareRoot{4608}}&\beta_{22}&=-\frac{\mathrm{J}_2+\mathrm{J}_6}{48}-\frac{\mathrm{J}_4}{\BareRoot{768}}+\frac{\mathrm{J}_1+\mathrm{J}_5}{\BareRoot{1536}}-\frac{\mathrm{J}_1-\mathrm{J}_5-2\mathrm{J}_3}{\BareRoot{4608}}\VthreeRaisedPoint
\end{aligned}
\]
\clearpage
\endgroup
\end{landscape}
\clearpage
\begin{landscape}
\begingroup
\fontsize{9.5}{12}\selectfont
\fancyhead[L]{\small NACHLASS. STÖRUNGEN DER PALLAS}
\renewcommand{\headrulewidth}{0pt}
\makeatletter\tagsleft@true\let\veqno\@@leqno\makeatother
\setlength{\lineskip}{4pt}\setlength{\lineskiplimit}{4pt}
\newcommand{\VthreeRaisedPoint}{\mathord{\text{\raisebox{.5ex}{.}}}}
\newcommand{\VthreeDZeroMark}{\mathord{\text{\makebox[0pt][l]{\kern-.2em\raisebox{1.25ex}{\rotatebox{10}{\rule{.04em}{.24ex}}}}}}}
\par\noindent\hspace*{\parindent}Weiter aus 4:
\[
\tag*{IV.}
\begin{aligned}
\tfrac12(\alpha_3+\alpha_{21})&=\frac{D_0-2D_8}{48}+\frac{D_2-D_6-D_{10}}{\BareRoot{1152}}&\;\tfrac12(\beta_3-\beta_{21})&=\frac{2E_4-E_{12}}{48}+\frac{E_2+E_6-E_{10}}{\BareRoot{1152}}\\[6pt]
\tfrac12(\alpha_9+\alpha_{15})&=\frac{D_0-2D_8}{48}-\frac{D_2-D_6-D_{10}}{\BareRoot{1152}}&\tfrac12(\beta_9-\beta_{15})&=-\frac{2E_4-E_{12}}{48}+\frac{E_2+E_6-E_{10}}{\BareRoot{1152}}\\[6pt]
\tfrac12(\alpha_1+\alpha_{23})&=\frac{D_0+D_8}{48}+\frac{D_4}{\BareRoot{768}}+\frac{D_2+D_{10}}{\BareRoot{1536}}+\frac{D_2-D_{10}+2D_6}{\BareRoot{4608}}&\tfrac12(\beta_1-\beta_{23})&=\frac{E_4+E_{12}}{48}+\frac{E_8}{\BareRoot{768}}+\frac{E_2+E_{10}}{\BareRoot{1536}}-\frac{E_2-E_{10}-2E_6}{\BareRoot{4608}}\\[6pt]
\tfrac12(\alpha_7+\alpha_{17})&=\frac{D_0+D_8}{48}-\frac{D_4}{\BareRoot{768}}-\frac{D_2+D_{10}}{\BareRoot{1536}}+\frac{D_2-D_{10}+2D_6}{\BareRoot{4608}}&\tfrac12(\beta_7-\beta_{17})&=-\frac{E_4+E_{12}}{48}+\frac{E_8}{\BareRoot{768}}+\frac{E_2+E_{10}}{\BareRoot{1536}}+\frac{E_2-E_{10}-2E_6}{\BareRoot{4608}}\\[6pt]
\tfrac12(\alpha_5+\alpha_{19})&=\frac{D_0+D_8}{48}-\frac{D_4}{\BareRoot{768}}+\frac{D_2+D_{10}}{\BareRoot{1536}}-\frac{D_2-D_{10}+2D_6}{\BareRoot{4608}}&\tfrac12(\beta_5-\beta_{19})&=\frac{E_4+E_{12}}{48}-\frac{E_8}{\BareRoot{768}}+\frac{E_2+E_{10}}{\BareRoot{1536}}+\frac{E_2-E_{10}-2E_6}{\BareRoot{4608}}\\[6pt]
\tfrac12(\alpha_{11}+\alpha_{13})&=\frac{D_0\VthreeDZeroMark+D_8}{48}+\frac{D_4}{\BareRoot{768}}-\frac{D_2+D_{10}}{\BareRoot{1536}}-\frac{D_2-D_{10}+2D_6}{\BareRoot{4608}}&\tfrac12(\beta_{11}-\beta_{13})&=-\frac{E_4+E_{12}}{48}-\frac{E_8}{\BareRoot{768}}+\frac{E_2+E_{10}}{\BareRoot{1536}}-\frac{E_2-E_{10}-2E_6}{\BareRoot{4608}}\VthreeRaisedPoint
\end{aligned}
\]
\clearpage
\endgroup
\end{landscape}
\clearpage
\begin{landscape}
\begingroup
\fontsize{9.5}{12}\selectfont
\fancyhead[L]{\small NACHLASS. STÖRUNGEN DER PALLAS}
\renewcommand{\headrulewidth}{0pt}
\makeatletter\tagsleft@true\let\veqno\@@leqno\makeatother
\setlength{\lineskip}{4pt}\setlength{\lineskiplimit}{4pt}
\newcommand{\VthreeRaisedPoint}{\mathord{\text{\raisebox{.5ex}{.}}}}
\newcommand{\VthreeUpperMark}{\mathord{\text{\makebox[.15em][l]{\kern.25em\raisebox{1.7ex}{\rule{.11em}{.055ex}}}}}}
\newcommand{\VthreeFractionLowPoint}{\mathord{\text{\makebox[0pt][l]{\kern-.25em\raisebox{-.1ex}{\rule{.07em}{.05ex}}}}}}
\par\noindent\hspace*{\parindent}Die Gleichungen 5 endlich lassen sich in folgender Weise auf eine für die numerische Rechnung geeignete Form bringen; es ist nemlich
\[
\begin{aligned}
\cos a+\cos7a&=\BareRoot{3}\mathord{.}\cos3a=\BareRoot{3}\mathord{.}\cos22^0 30'&\qquad\sin a+\sin7a&=\sin9a=\sin67^0 30'\\[4pt]
\cos a-\cos7a&=\cos9a=\cos67^0 30'&\sin a-\sin7a&=-\BareRoot{3}\mathord{.}\sin3a=-\BareRoot{3}\mathord{.}\sin22^0 30'\\[4pt]
\cos5a+\cos11a&=\sin9a=\sin67^0 30'&\sin5a+\sin11a&=\BareRoot{3}\mathord{.}\cos3a=\BareRoot{3}\mathord{.}\cos22^0 30'\\[4pt]
\cos5a-\cos11a&=\BareRoot{3}\mathord{.}\sin3a=\BareRoot{3}\mathord{.}\sin22^0 30'&\sin5a-\sin11a&=-\cos9a=-\cos67^0 30'.
\end{aligned}
\]
Setzt man also:
\SourceLeaf{505}
\[
\begin{aligned}
D_1+D_7&=a\cos A&\qquad-E_1+E_7&=a'\sin A'\\
D_5-D_{11}&=a\sin A&E_5+E_{11}&=a'\cos A'\\[6pt]
D_1-D_7-D_9&=b\cos B&E_1+E_7-E_9&=b'\sin B'\\
D_3-D_5-D_{11}&=b\sin B&E_3+E_5-E_{11}&=b'\cos B'\\[6pt]
D_1-D_7+2D_9&=c\cos C&E_1+E_7+2E_9&=c'\sin C'\\
2D_3+D_5+D_{11}&=c\sin C&2E_3-E_5+E_{11}&=c'\cos C',
\end{aligned}
\]
so überzeugt man sich leicht, dass die letztgenannten Formeln in die folgenden übergehen:
\[
\tag*{V.}
\begin{aligned}
24\mathord{.}\tfrac12(\alpha_3-\alpha_{21})&=b\cos(22^0 30'-B)&\qquad24\mathord{.}\tfrac12(\beta_3+\beta_{21})&=b'\cos(B'-22^0 30')\\[4pt]
24\mathord{.}\tfrac12(\alpha_9-\alpha_{15})&=b\sin(22^0 30'-B)&24\mathord{.}\tfrac12(\beta_9+\beta_{15})&=b'\sin(B'-22^0 30')\\[6pt]
\BareRoot{768}\mathord{.}\tfrac14(\alpha_1+\alpha_7-\alpha_{17}-\alpha_{23})&=a\cos(22^0 30'-A)&\BareRoot{768}\mathord{.}\tfrac14(\beta_1-\beta_7-\beta_{17}+\beta_{23})&=a'\cos(22^0 30'-A')\\[4pt]
\BareRoot{768}\mathord{.}\tfrac14(\alpha_5-\alpha_{11}+\alpha_{13}-\alpha_{19})&=a\sin(22^0 30'-A)&\BareRoot{768}\mathord{.}\tfrac14(\beta_5+\beta_{11}+\beta_{13}+\beta_{19})&=a'\sin(22^0 30'-A')\\[6pt]
48\mathord{.}\tfrac14(\alpha_1-\alpha_7+\alpha_{17}-\alpha_{23})&=c\cos(67^0 30'-C)&48\mathord{.}\tfrac14(\beta_1+\beta_7+\beta_{17}+\beta_{23})&=c'\cos(67^0 30'-C')\\[4pt]
48\mathord{.}\tfrac14(\alpha_5+\alpha_{11}-\alpha_{13}-\alpha_{19})&=c\sin(67^0 30'-C)&48\mathord{.}\tfrac14(\beta_5-\beta_{11}-\beta_{13}+\beta_{19})&=c'\sin(67^0 30'-C').
\end{aligned}
\]
\par\noindent\hspace*{\parindent}Für \(\mu=24\), also \(a=15^0\), lassen sich in ganz analoger Weise die folgenden Formeln ableiten:
\[
\begin{aligned}
\alpha_6&=\frac{H_0-2H_2}{24}&\qquad\beta_6&=\frac{2\mathrm{J}_1-\mathrm{J}_3}{24}\VthreeUpperMark\\[6pt]
\alpha_2&=\frac{H_0+H_2}{24}+\frac{H_1}{\BareRoot{192}}&\beta_2&=\frac{\mathrm{J}_1+\mathrm{J}_3}{24}+\frac{\mathrm{J}_2}{\BareRoot{192}}\\[6pt]
\alpha_{10}&=\frac{H_0+H_2}{24}-\frac{H_1}{\BareRoot{192}}&\beta_{10}&=\frac{\mathrm{J}_1+\mathrm{J}_3}{24}-\frac{\mathrm{J}_2}{\BareRoot{192}}\\[6pt]
\alpha_0&=\frac{P_0+2P_1}{48}&&\\[6pt]
\alpha_4&=\frac{M_0+M_1}{24}&\beta_4&=\frac{N_1}{\BareRoot{192}}\\[6pt]
\alpha_8&=\frac{P_0-P_1}{24}&\beta_8&=\frac{Q_1}{\BareRoot{192}}\\[6pt]
\alpha_{12}&=\frac{M_0-2M_1}{48}&&
\end{aligned}
\]
\clearpage
\[
\begin{aligned}
\tfrac12(\alpha_3+\alpha_9)&=\frac{D_0-2D_4}{24}&\tfrac12(\beta_3-\beta_9)&=\frac{2E_2-E_6}{24}\\[6pt]
\tfrac12(\alpha_3-\alpha_9)&=\frac{D_1-D_3-D_5}{\BareRoot{288}}&\tfrac12(\beta_3+\beta_9)&=\frac{E_1+E_3-E_5}{\BareRoot{288}}\\[10pt]
\tfrac12(\alpha_1+\alpha_{11})&=\frac{D_0+D_4}{24}+\frac{D_2}{\BareRoot{192}}&\tfrac12(\beta_1-\beta_{11})&=\frac{E_2+E_6}{24}+\frac{E_4}{\BareRoot{192}}\\[6pt]
\tfrac12(\alpha_5+\alpha_7)&=\frac{D_0+D_4}{24}-\frac{D_2}{\BareRoot{192}}&\tfrac12(\beta_5-\beta_7)&=\frac{E_2+E_6}{24}-\frac{E_4}{\BareRoot{192}}\\[6pt]
\tfrac12(\alpha_1-\alpha_{11})&=\frac{D_1+D_5}{\BareRoot{384}}+\frac{D_1-D_5+2D_3}{\BareRoot{1152}}&\tfrac12(\beta_1+\beta_{11})&=\frac{E_1+E_5}{\BareRoot{384}}-\frac{E_1-E_5-2E_3}{\BareRoot{1152}}\\[6pt]
\tfrac12(\alpha_5-\alpha_7)&=\frac{D_1+D_5}{\BareRoot{384}}-\frac{D_1-D_5+2D_3}{\BareRoot{1152}}&\tfrac12(\beta_5+\beta_7)&=\frac{E_1+E_5}{\BareRoot{384}}+\frac{E_1-E_5-2E_3}{\BareRoot{1152}}\VthreeRaisedPoint
\end{aligned}
\]
\clearpage\SourceLeaf{506}
\par\noindent\hspace*{\parindent}Für \(\mu=12\), also \(a=30^0\), endlich wird:
\[
\begin{aligned}
\alpha_2&=\frac{H_0+H_1}{12}&\qquad\beta_2&=\frac{\mathrm{J}_1}{\BareRoot{48}}\\[6pt]
\alpha_6&=\frac{H_0-2H_1}{24}&&\\[6pt]
\alpha_4&=\frac{K_0-K_1}{12}&\beta_4&=\frac{L_1}{\BareRoot{48}}\\[6pt]
\alpha_0&=\frac{K_0+2K_1}{24}&&\\[6pt]
\alpha_3&=\frac{D_0-2D_2}{12}&\beta_3&=\frac{2E_1-E_3}{12}\VthreeFractionLowPoint\\[6pt]
\tfrac12(\alpha_1+\alpha_5)&=\frac{D_0+D_2}{12}&\tfrac12(\beta_1+\beta_5)&=\frac{E_1+E_3}{12}\\[6pt]
\tfrac12(\alpha_1-\alpha_5)&=\frac{D_1}{\BareRoot{48}}&\tfrac12(\beta_1-\beta_5)&=\frac{E_2}{\BareRoot{48}}\VthreeRaisedPoint\Bigg]
\end{aligned}
\]
\clearpage
\endgroup
\end{landscape}
